\kefalaio{Ευκλείδεια διαίρεση: το υπόλοιπο που μένει}{Κεφάλαιο Α.1.4, «Ευκλείδεια διαίρεση – Διαιρετότητα»}

\begin{anagke}
Το σχολείο οργανώνει την πρώτη του εκδρομή. Θα πάνε και τα 387 παιδιά, με λεωφορεία των 49 θέσεων. Η κυρία Αλεξάνδρα ζητάει από την τάξη να βρει πόσα λεωφορεία χρειάζονται.

Ο Στέλιος, που αγαπάει το κομπιουτεράκι, κάνει $387:49$ και η οθόνη γράφει $7{,}897959\dots$. «Εφτά κόμμα εννιά!» λέει περήφανος. «Άρα εφτά λεωφορεία». Η Μυρτώ τον κοιτάει. «Και τα υπόλοιπα παιδιά; Μένουν στην αυλή;» Ο Στέλιος ξαναδιαβάζει την οθόνη. «Μα το 0,897 τι είναι; Δεν υπάρχει 0,897 λεωφορείου». Η Ζωή κάνει τη διαίρεση στο χαρτί, όπως τη μάθαμε: πηλίκο 7 και κάτι περισσεύει. «Αυτό που περισσεύει είναι παιδιά, όχι κόμματα. Εκεί βρίσκεται η απάντηση».

Σε μια διαίρεση, πολλές φορές, δεν ενδιαφέρει μόνο πόσες φορές χωράει ο ένας αριθμός στον άλλον. Ενδιαφέρει και τι περισσεύει. Πώς γράφεται αυτό με κανόνα;
\end{anagke}

\begin{lexi}{πηλίκο, υπόλοιπο}{πηλίκο, υπόλοιπο}
Το \lexiword{πηλίκο} έρχεται από το αρχαίο «πηλίκος», που σημαίνει «πόσο μεγάλος». Η διαίρεση απαντάει στο «πόσο μεγάλο μερίδιο παίρνει καθένας». Το \lexiword{υπόλοιπο} είναι κυριολεκτικά αυτό που μένει: υπό-λοιπο, το «λοιπό» που δεν μοιράστηκε.

Η \lexiword{Ευκλείδεια διαίρεση} πήρε το όνομά της από τον Ευκλείδη, τον Έλληνα μαθηματικό που έζησε στην Αλεξάνδρεια περίπου το 300 π.Χ. και έγραψε τα «Στοιχεία», το βιβλίο όπου συγκέντρωσε όλα τα μαθηματικά της εποχής του. Η Ευκλείδεια διαίρεση είναι η διαίρεση με πηλίκο και υπόλοιπο, όπως τη μαθαίνουμε στο δημοτικό.
\end{lexi}

\begin{metafrasi}
Μοιράζω 43 μπάλες σε ομάδες των 7: γίνονται 6 ομάδες και μένει 1 μπάλα. & $43=7\cdot 6+1$\\
Δεν περισσεύει τίποτα (τέλεια διαίρεση). & $\Delta=\delta\cdot\pi$, υπόλοιπο 0\\
Το υπόλοιπο είναι πάντα μικρότερο από τον διαιρέτη. & $\upsilon<\delta$\\
Στη διαίρεση με το 8 τα πιθανά υπόλοιπα. & $0,\,1,2,\,3,4,\,5,6,\,7$\\
Δεν διαιρούμε ποτέ με το μηδέν. & $\delta\neq 0$\\
\end{metafrasi}

\tmima{Η θεωρία}

\begin{orismos}{Ευκλείδεια διαίρεση}
Όταν δίνονται δύο φυσικοί $\Delta$ και $\delta$ (με $\delta\neq 0$), υπάρχουν δύο φυσικοί $\pi$ και $\upsilon$ με
$$\Delta=\delta\cdot\pi+\upsilon,\qquad \upsilon<\delta.$$
Ο $\Delta$ λέγεται \textbf{διαιρετέος}, ο $\delta$ \textbf{διαιρέτης}, ο $\pi$ \textbf{πηλίκο} και ο $\upsilon$ \textbf{υπόλοιπο}.
\end{orismos}

Η συνθήκη $\upsilon<\delta$ είναι το κλειδί. Αν το υπόλοιπο ήταν ίσο ή μεγαλύτερο από τον διαιρέτη, θα χωρούσε άλλη μία φορά ο διαιρέτης και θα μπορούσαμε να αυξήσουμε το πηλίκο. Η \textbf{δοκιμή} της διαίρεσης είναι η ίδια η ισότητα: $\delta\cdot\pi+\upsilon$ πρέπει να δίνει τον $\Delta$.

Όταν $\upsilon=0$ έχουμε \textbf{τέλεια διαίρεση}: $\Delta=\delta\cdot\pi$. Τότε λέμε ότι ο $\Delta$ \textbf{διαιρείται} με τον $\delta$, και ότι ο $\delta$ \textbf{διαιρεί} τον $\Delta$. Στην τέλεια διαίρεση η διαίρεση είναι η αντίστροφη πράξη του πολλαπλασιασμού: αν $\Delta=\delta\cdot\pi$, τότε $\Delta:\delta=\pi$ και $\Delta:\pi=\delta$.

\begin{idiotita}{διαίρεση}
\begin{itemize}[leftmargin=1.2em,itemsep=0.2ex,topsep=2pt]
\item Ο διαιρέτης δεν μπορεί να είναι 0: δεν υπάρχει $\Delta:0$.
\item $\alpha:\alpha=1$ (αν $\alpha\neq 0$). \quad $\alpha:1=\alpha$. \quad $0:\alpha=0$ (αν $\alpha\neq 0$).
\end{itemize}
\end{idiotita}

Γιατί δεν διαιρούμε με το 0; Αν ήταν $6:0=x$, τότε θα έπρεπε $6=0\cdot x$. Μα κάθε αριθμός επί 0 κάνει 0, ποτέ 6. Δεν υπάρχει κανένας τέτοιος $x$.

\begin{prosoxi}
\begin{itemize}[leftmargin=1.2em,itemsep=0.3ex,topsep=2pt]
\item Η ισότητα $\Delta=\delta\cdot\pi+\upsilon$ δεν αρκεί για να είναι Ευκλείδεια διαίρεση. Πρέπει και $\upsilon<\delta$. Το $1.345=59\cdot 21+106$ είναι σωστή ισότητα, αλλά το 106 είναι μεγαλύτερο από το 59 και από το 21: δεν είναι διαίρεση.
\item Το πηλίκο δεν είναι πάντα η απάντηση σε ένα πρόβλημα. Για 387 παιδιά και λεωφορεία των 49 θέσεων, το πηλίκο είναι 7, αλλά χρειάζονται 8 λεωφορεία, γιατί το υπόλοιπο παιδιά θέλει θέσεις. Διάβασε τι ρωτάει το πρόβλημα.
\end{itemize}
\end{prosoxi}

\tmima{Λυμένα παραδείγματα}

\paradeigma{Τα λεωφορεία}
\tixero 387 παιδιά, λεωφορεία με 49 θέσεις.\par
\tithelo Πόσα λεωφορεία χρειάζονται και πόσες θέσεις θα μείνουν άδειες.\par
\possindeo Κάνω την Ευκλείδεια διαίρεση $387:49$. Το $49\cdot 7=343$ και το $49\cdot 8=392$ είναι πάνω από το 387, άρα πηλίκο 7 και υπόλοιπο $387-343=44$. Γράφω $387=49\cdot 7+44$. Επτά λεωφορεία γεμίζουν, και 44 παιδιά περισσεύουν. Χρειάζεται ένα όγδοο λεωφορείο γι' αυτά. Σε αυτό θα μείνουν $49-44=5$ θέσεις άδειες.\par
\apantisi 8 λεωφορεία, και 5 άδειες θέσεις.\par
\elegxos Η δοκιμή: $49\cdot 7+44=343+44=387$. Και $44<49$, όπως πρέπει.

\paradeigma{Ποιες ισότητες είναι διαιρέσεις;}
\tixero Οι ισότητες α) $120=28\cdot 4+8$, β) $1.345=59\cdot 21+106$, γ) $374=8\cdot 46+6$.\par
\tithelo Σε ποιες περιπτώσεις είναι Ευκλείδεια διαίρεση και με ποιον διαιρέτη.\par
\possindeo Ελέγχω ότι η ισότητα ισχύει και ότι το υπόλοιπο είναι μικρότερο από τον διαιρέτη.\par
α) Το υπόλοιπο 8 είναι μικρότερο από το 28, όχι από το 4. Είναι διαίρεση με διαιρέτη το 28 και πηλίκο 4.\par
β) Το υπόλοιπο 106 είναι μεγαλύτερο και από το 59 και από το 21. Δεν είναι διαίρεση.\par
γ) Το υπόλοιπο 6 είναι μικρότερο και από το 8 και από το 46. Είναι διαίρεση και με διαιρέτη το 8 (πηλίκο 46) και με διαιρέτη το 46 (πηλίκο 8).\par
\apantisi Μόνο οι (α) και (γ).\par
\elegxos Για το (β): $59\cdot 21=1.239$ και $1.345-1.239=106$. Ναι, η ισότητα ισχύει, αλλά το 106 περισσεύει πολύ: ο 59 χωράει ακόμα μία φορά στο 106. Άρα το πηλίκο δεν ήταν 21.

\paradeigma{Η μέρα της εβδομάδας}
\tixero Σήμερα είναι Τρίτη.\par
\tithelo Τι μέρα θα είναι μετά από 247 ημέρες.\par
\possindeo Η εβδομάδα έχει 7 μέρες, και κάθε 7 μέρες γυρίζουμε στην ίδια μέρα. Διαιρώ: $247=7\cdot 35+2$. Οι 35 ολόκληρες εβδομάδες δεν αλλάζουν τη μέρα. Μένουν 2 μέρες: Τετάρτη, Πέμπτη.\par
\apantisi Πέμπτη.\par
\elegxos $7\cdot 35=245$, και μετά από 245 ημέρες είναι πάλι Τρίτη. Δύο μέρες μετά είναι Πέμπτη. Το υπόλοιπο μας είπε πόσο προχωράμε μέσα στην εβδομάδα.

\tmima{Ασκήσεις}

\zoni{Το κατάλαβα}{εφαρμογή του κανόνα}

\ex Κάνε τις διαιρέσεις και τη δοκιμή τους: α) $4.002:69$ \quad β) $1.445:17$ \quad γ) $925:37$ \quad δ) $3.621:213$ \quad ε) $35.280:2.940$ \quad στ) $5.082:77$.

\ex Βρες πηλίκο και υπόλοιπο και γράψε την ισότητα $\Delta=\delta\cdot\pi+\upsilon$: α) $43:7$ \quad β) $100:9$ \quad γ) $245:12$ \quad δ) $1.000:33$ \quad ε) $59:60$ \quad στ) $0:5$.

\ex Ποιες από τις παρακάτω ισότητες εκφράζουν Ευκλείδεια διαίρεση; Με ποιον διαιρέτη; α) $125=35\cdot 3+20$ \quad β) $762=38\cdot 19+40$ \quad γ) $1.500=42\cdot 35+30$ \quad δ) $300=18\cdot 16+12$.

\ex Αν $\nu$ είναι ένας φυσικός αριθμός, ποια υπόλοιπα μπορεί να έχει η διαίρεση $\nu:8$;

\ex Ένας αριθμός διαιρούμενος με το 9 δίνει πηλίκο 73 και υπόλοιπο 4. Ποιος είναι ο αριθμός;

\ex Συμπλήρωσε: α) $\delta=12$, $\pi=7$, $\upsilon=5$: βρες τον $\Delta$. \quad β) $\Delta=100$, $\delta=8$: βρες $\pi$ και $\upsilon$. \quad γ) $\Delta=59$, $\pi=4$, $\upsilon=3$: βρες τον $\delta$.

\ex Σωστό ή λάθος; α) Στη διαίρεση με το 6 το υπόλοιπο μπορεί να είναι 6. \quad β) $0:7=0$. \quad γ) $7:0=0$. \quad δ) $\alpha:\alpha=1$ για κάθε $\alpha\neq 0$. \quad ε) $17:1=1$. \quad στ) Όταν το υπόλοιπο είναι 0, η διαίρεση λέγεται τέλεια.

\ex Ο Στέλιος γράφει $85=12\cdot 6+13$ και λέει «είναι διαίρεση, γιατί η ισότητα ισχύει». Ισχύει η ισότητα; Γιατί παρ' όλα αυτά δεν είναι Ευκλείδεια διαίρεση; Διόρθωσέ την.

\zoni{Το χρησιμοποιώ}{προβλήματα από τη ζωή}

\ex \textbf{Η εκδρομή.} Στην εκδρομή του σχολείου πάνε 387 παιδιά και κάθε λεωφορείο έχει 49 θέσεις. Πόσα λεωφορεία χρειάζονται; Πόσες θέσεις μένουν κενές στο τελευταίο; Αν η εταιρεία έχει μόνο 7 λεωφορεία, πόσα παιδιά θα μείνουν στο σχολείο;

\ex \textbf{Οι μέρες.} Σήμερα είναι Τρίτη. Τι μέρα θα είναι μετά από 150 ημέρες; Τι μέρα ήταν πριν από 8 ημέρες;

\ex \textbf{Οι γλάστρες του παππού.} Ο παππούς έχει 1.000 σπόρους και φυτεύει 33 σε κάθε γλάστρα. Πόσες γλάστρες γεμίζει; Πόσοι σπόροι περισσεύουν; Φτάνουν οι περισσευούμενοι για μία ακόμα γλάστρα;

\ex \textbf{Οι μπάλες.} Ο Άρης έχει 87 € και θέλει να αγοράσει μπάλες του μπάσκετ που κοστίζουν 13 € η καθεμία. Πόσες μπάλες αγοράζει και πόσα ευρώ του μένουν;

\ex \textbf{Η παράταξη.} Στην αυλή θα παραταχθούν 168 μαθητές της Α’ τάξης για μια φωτογραφία. Ο γυμναστής ρωτάει αν μπορούν να σταθούν σε ίσες ομάδες των 3, των 4, των 5, των 6 ή των 7, χωρίς να περισσεύει κανείς. Κάνε τις διαιρέσεις και απάντησε για κάθε περίπτωση.

\ex \textbf{Το ρολόι.} Η Μυρτώ βάζει το ξυπνητήρι στις 10:00 το πρωί. Θέλει να δει τι ώρα θα είναι μετά από 50 ώρες. Σκέψου ότι κάθε 24 ώρες το ρολόι γυρίζει στην ίδια ώρα.

\zoni{Το σκέφτομαι}{μια ερώτηση που πάει πιο μακριά}

\ex Βρες όλους τους φυσικούς μεταξύ του 20 και του 40 που, διαιρούμενοι με το 7, αφήνουν υπόλοιπο 3. Τι παρατηρείς για τη διαφορά ανάμεσα σε δύο διαδοχικούς από αυτούς;

\ex Ένας αριθμός $\nu$ αφήνει υπόλοιπο 4 όταν διαιρεθεί με το 6. Δοκίμασε μερικές τιμές του $\nu$ (10, 16, 22, \dots). Τι υπόλοιπο αφήνει ο $\nu$ όταν διαιρεθεί με το 3;

\ex Στη διαίρεση ενός αριθμού με έναν άλλον το πηλίκο είναι 12 και το υπόλοιπο 7. Ποιος είναι ο μικρότερος δυνατός διαιρέτης; Ποιος είναι τότε ο διαιρετέος;

\ex Η Ζωή ρωτάει: «Γιατί δεν διαιρούμε με το 0;» Ο Ηλίας απαντάει «γιατί δεν το επιτρέπει ο δάσκαλος». Εξήγησε με το $6:0=x$ και την αντίστροφη πράξη ότι δεν είναι θέμα άδειας. Τι γίνεται στο $0:0$;

\ex Σήμερα είναι Τρίτη. Τι μέρα θα είναι μετά από ένα εκατομμύριο ημέρες; (Μην αρχίσεις να μετράς. Βρες το υπόλοιπο της διαίρεσης του $1.000.000$ με το 7.)
